% Generated by scripts/assemble.py; corrections belong in source/corrections.json.
\examhead{86}{东南大学 2024 年数学分析试卷}{84}
% p085-b19; source PDF page 85
\par\Needspace{5\baselineskip}
\blockanchor{p085-b19}
\majorquestion{一}计算题与解答题. 每题 9 分, 共 90 分.\par
% p085-b20; source PDF page 85
\par\Needspace{6\baselineskip}
\blockanchor{p085-b20}
\question{1}设 $y = y(x)$ 是由方程\par
\nopagebreak[4]
% p085-b21; source PDF page 85
\blockanchor{p085-b21}
\displayformula{\arcsin (x y) + y ^ {2} - y e ^ {x y} = 2}
% p085-b22; source PDF page 85
\blockanchor{p085-b22}
所确定的隐函数, 求曲线 $y = y(x)$ 在点 (0,2) 处的切线方程.\par
% p085-b23; source PDF page 85
\blockanchor{p085-b23}
\question{2}求极限 $\lim_{n\to \infty}\left(\sqrt[n]{n^2 + 1} -1\right)\sin \frac{n\pi}{2}$.\par
% p085-b24; source PDF page 85
\blockanchor{p085-b24}
\question{3}计算不定积分 $\int \frac{\arccos x}{x^2} \, \mathrm{d}x$.\par
% p085-b25; source PDF page 85
\blockanchor{p085-b25}
\question{4}计算定积分 $\int_{-1}^{1} \frac{x + 2}{e^x + e^{-x}} \, \mathrm{d}x$.\par
% p086-b00; source PDF page 86
\blockanchor{p086-b00}
\question{5}求幂级数 $\sum_{n=1}^{\infty} \frac{n+1}{n!2^n}(x-2)^n$ 的收敛域与和函数.\par
% p086-b01; source PDF page 86
\blockanchor{p086-b01}
\question{6}计算二重积分 $\iint_{D} \frac{\sin y \cos y}{y} \, \mathrm{d}x \, \mathrm{d}y$, 其中 $D$ 为直线 $y = x$ 与抛物线 $x = y^2$ 所围成的区域.\par
% p086-b02; source PDF page 86
\blockanchor{p086-b02}
\question{7}设 $f(x, y) = \begin{cases} \frac{x^2y}{\sqrt{x^2 + y^2}}, & x^2 + y^2 \neq 0; \\ 0, & x^2 + y^2 = 0. \end{cases}$ 求二阶混合偏导数 $f_{xy}$.\par
% p086-b03; source PDF page 86
\blockanchor{p086-b03}
\question{8}设 $f(x) = \pi - x, x \in [0, \pi]$.\par
% p086-b04; source PDF page 86
\blockanchor{p086-b04}
(1) 将 $f(x)$ 展开为余弦级数, 并在 $[- \pi, \pi]$ 上写出和函数的表达式.\par
% p086-b05; source PDF page 86
\blockanchor{p086-b05}
(2) 判断该级数在 $[0, \pi]$ 内是否一致收敛，并说明原因.\par
% p086-b06; source PDF page 86
\par\Needspace{6\baselineskip}
\blockanchor{p086-b06}
\question{9}计算曲线积分\par
\nopagebreak[4]
% p086-b07; source PDF page 86
\blockanchor{p086-b07}
\displayformula{\int_ {L} y ^ {2} \mathrm{d} x + z ^ {2} \mathrm{d} y + x ^ {2} \mathrm{d} z.}
% p086-b08; source PDF page 86
\blockanchor{p086-b08}
其中 $L$ 为球面 $x^{2} + y^{2} + z^{2} = a^{2}$ ($z \geq 0, a > 0$) 与柱面 $x^{2} + y^{2} - ax = 0$ 的交线, 从 $z$ 轴正向看去取逆时针方向.\par
% p086-b09; source PDF page 86
\blockanchor{p086-b09}
\question{10}解答如下问题:\par
% p086-b10; source PDF page 86
\blockanchor{p086-b10}
(1) 求极限  $\lim_{x\to0^{+}}\frac{1-\cos x}{\int_{0}^{x}\frac{\ln(1+xy)}{y}dy}.$\par
% p086-b11; source PDF page 86
\blockanchor{p086-b11}
(2) 计算含参量反常积分  $\int_{0}^{+\infty}\frac{\sin xy}{ye^{y}}dy$ .\par
% p086-b12; source PDF page 86
\par\Needspace{5\baselineskip}
\blockanchor{p086-b12}
\majorquestion{二}证明题. 每题 10 分, 共 50 分.\par
% p086-b13; source PDF page 86
\par\Needspace{6\baselineskip}
\blockanchor{p086-b13}
\question{11}设 $f(x)$ 在 $[0,1]$ 上连续，在 $(0,1)$ 内可导，\par
\nopagebreak[4]
% p086-b14; source PDF page 86
\blockanchor{p086-b14}
\displayformula{f (0) = f (1) = 0, f \left(\frac {1}{2}\right) = 1.}
% p086-b15; source PDF page 86
\blockanchor{p086-b15}
证明：存在 $\xi \in (0,1)$，使得 $f'(\xi) = 1$\par
% p086-b16; source PDF page 86
\par\Needspace{6\baselineskip}
\blockanchor{p086-b16}
\question{12}设 $D = \{(x, y) \mid 0 < x < 1, 0 < y < +\infty\}$, 证明: 对任意的 $(x, y) \in D$, 成立不等式\par
\nopagebreak[4]
% p086-b17; source PDF page 86
\blockanchor{p086-b17}
\displayformula{y x ^ {y} (1 - x) <   e ^ {- 1}.}
% p086-b18; source PDF page 86
\blockanchor{p086-b18}
\question{13}设 $f_{n}(x) = n^{\alpha}xe^{-nx}$.\par
% p086-b19; source PDF page 86
\blockanchor{p086-b19}
(1) $\alpha$ 取何值时, $\lim_{n\to \infty}\int_0^1 f_n(x)\mathrm{d}x = \int_0^1\lim_{n\to \infty}f_n(x)\mathrm{d}x?$\par
% p086-b20; source PDF page 86
\blockanchor{p086-b20}
(2) $\alpha$ 取何值时, $\{f_n(x)\}$ 在[0,1]上一致收敛?\par
% p086-b21; source PDF page 86
\blockanchor{p086-b21}
\question{14}设正项级数 $\sum_{n=1}^{\infty} a_n$ 发散, 以 $S_n$ 表示前 $n$ 项的和. 证明:\par
% p086-b22; source PDF page 86
\blockanchor{p086-b22}
(1) $\sum_{n=1}^{\infty} \frac{a_n}{S_n}$ 发散.\par
% p086-b23; source PDF page 86
\blockanchor{p086-b23}
(2) $\sum_{n=1}^{\infty} \frac{a_n}{S_n^2}$ 收敛.\par
% p086-b24; source PDF page 86
\blockanchor{p086-b24}
\question{15}设 $f(x, y)$ 在区域 $D: (0, 1) \times (0, 1)$ 内分别对自变量 $x$ 和 $y$ 连续，并且对固定的 $x, f(x, y)$ 关于 $y$ 是单调的。证明：$f(x, y)$ 在区域 $D$ 内为二元连续函数。\par
% p086-b25; source PDF page 86
\blockanchor{p086-b25}
\question{16}解答如下问题:\par
% p086-b26; source PDF page 86
\blockanchor{p086-b26}
(1) 叙述  $R^{n}$  上的有限覆盖定理.\par
% p086-b27; source PDF page 86
\blockanchor{p086-b27}
(2) 设对任意的 $x_0 \in [a, b]$, 有 $\lim_{x \to x_0} f(x) = 0$. 证明: $f(x) \in R[a, b]$ 且 $\int_{a}^{b} f(x) \mathrm{d}x = 0$.\par
